\chapter{Command-Line Reference}
\label{app:cli}

\section{Building}

\begin{center}
\small
\begin{tabular}{@{}lp{0.60\textwidth}@{}}
\toprule
\ctp{make} & build \ctp{./st80}\\
\ctp{make test} & unit suites, then every profile's own SUnit suites\\
\ctp{make unit-test} & just the unit suites\\
\ctp{make suite-test} & just the imported packages' suites\\
\ctp{make bench} & the parallel scaling benchmark\\
\ctp{make font} & regenerate the built-in face (needs Pillow and a font)\\
\ctp{make clean} & remove build artifacts\\
\ctp{make help} & targets and variables\\
\bottomrule
\end{tabular}
\end{center}

\begin{center}
\small
\begin{tabular}{@{}lp{0.60\textwidth}@{}}
\toprule
\ctp{OM=mt} & 64-bit threaded object memory --- the system, and the default\\
\ctp{OM=bb} & 16-bit Blue Book memory --- the Xerox trace harness. Cannot
  bootstrap an image\\
\ctp{TSAN=1} & ThreadSanitizer build\\
\ctp{ASAN=1} & Address and undefined-behaviour sanitizers\\
\ctp{OPT=-O0} & override optimisation flags\\
\ctp{FONT=} \ctp{SIZE=} \ctp{LEAD=} & inputs to \ctp{make font}\\
\bottomrule
\end{tabular}
\end{center}

\section{\texttt{st80}}

\begin{center}
\small
\begin{tabular}{@{}lp{0.56\textwidth}@{}}
\toprule
\ctp{-version} & version and build configuration\\
\ctp{-help} & the option list\\
\midrule
\ctp{-bootstrap <a.st...>} & build an image from source\\
\ctp{\ \ -profile p} & use a profile\\
\ctp{\ \ -manifest f} & use a 1983-style manifest\\
\ctp{\ \ -o image} & write the image (needs \ctp{OM=mt})\\
\ctp{\ \ -eval expr} & evaluate and print\\
\ctp{\ \ -startup expr} & what a saved image evaluates when it resumes\\
\ctp{\ \ -closures} & compile in the closures dialect\\
\ctp{\ \ -tests} & run every SUnit test; exit non-zero if any failed\\
\ctp{\ \ -doctests f.st} & run the \ctp{expr >>> value} examples in a file\\
\midrule
\ctp{-run <image> [n]} & run the image, opening a window. \ctp{n} is a bytecode limit, not a worker count\\
\ctp{-serve <image> [-workers n] [args]} & run the image on \ctp{n} native threads (one per CPU less one by default), no window, until SIGINT, SIGTERM or \ctp{Smalltalk quitPrimitive}; \ctp{args} are what \ctp{Smalltalk arguments} answers. The image's startup is what \ctp{-bootstrap -startup} gave it---Chapter~\ref{ch:serving}\\
\ctp{-screenshot f.pbm} & with \ctp{-run} or \ctp{-bootstrap}, write the
  display\\
\ctp{-inject <script>} & post input events\\
\ctp{-wiggle} & move the pointer continuously, no clicks\\
\midrule
\ctp{-syntax <f.st...>} & compile every method and report failures\\
\ctp{-primitives <f.st...>} & every primitive that source asks the VM for\\
\ctp{-census <image>} & load an image and summarise it\\
\ctp{-classes <image>} & every class, in \ctp{class.oops} format\\
\ctp{-methods <image>} & every method, in \ctp{method.oops} format\\
\ctp{-inspect <image> <oop>} & describe one object (oop in hex)\\
\ctp{-disasm <image> <Class> <sel>} & dump a method's bytecodes\\
\ctp{-trace2 <image> [n]} & bytecode trace, Xerox \ctp{trace2} format\\
\ctp{-trace3 <image> [n]} & send trace, Xerox \ctp{trace3} format\\
\bottomrule
\end{tabular}
\end{center}

\section{Input scripts}

For \ct{-inject}, semicolon-separated:

\begin{center}
\small
\begin{tabular}{@{}ll@{}}
\toprule
\ctp{m X Y} & move the pointer to X, Y\\
\ctp{d CODE} & press a button: 128 right, 129 middle, 130 left\\
\ctp{u CODE} & release it\\
\ctp{k CODE} & type a key\\
\ctp{w N} & wait N bytecode slices\\
\ctp{x} & post a window-exposed event\\
\bottomrule
\end{tabular}
\end{center}

\section{Environment}

\begin{center}
\small
\begin{tabular}{@{}lp{0.56\textwidth}@{}}
\toprule
\ctp{ST\_DISPLAY\_THEME} & \ctp{paper} (default), \ctp{classic}, \ctp{dark}\\
\ctp{ST\_DISPLAY\_SCALE} & screen pixels per display pixel\\
\ctp{ST\_DISPLAY\_WINDOW} & \ctp{WxH}: open the window at this size\\
\ctp{ST\_DISPLAY\_FIT=off} & keep the image's own screen size\\
\ctp{ST\_DISPLAY\_PRESENTATION} & \ctp{integer}, \ctp{letterbox},
  \ctp{stretch}\\
\ctp{ST\_DISPLAY\_TRACE=1} & report why a resize was refused; draw and present
  counts at exit\\
\ctp{ST\_EVAL\_TRACE=1} & trace the bytecodes an \ctp{-eval} runs\\
\ctp{ST\_NET\_TRACE=1} & narrate the network thread: arms, deliveries, closes\\
\ctp{ST\_BOTTOM\_LOG=1} & report a return off the bottom of a stack\\
\bottomrule
\end{tabular}
\end{center}
